\chapter{DataFrames}
\label{cap:dataframes}

\section{What is a DataFrame}

A DataFrame is a two-dimensional table with named columns. In \hayashi{},
each column is a 64-bit floating-point vector. Missing values are
represented by \texttt{NaN} (\textit{Not a Number}).

The DataFrame is the central type of the language — virtually every
econometric analysis starts with a DataFrame loaded from some source.

\section{Copy semantics (COW)}

DataFrames use copy-on-write (COW) semantics: assigning a DataFrame to
another variable creates an alias that shares the same internal data. The
actual copy occurs only when one of them is modified.

\begin{lstlisting}[caption={COW: baseline preserved}]
load "panel.csv" as df

let baseline = df      // cheap alias, no copy

let df2 = df |> generate(z = x^2)    // df2 has 'z', df does not
let df3 = df |> filter(year > 2000)  // df3 filtered, df intact

// baseline == original df at any point
\end{lstlisting}

This semantics allows creating analysis variants with no memory cost until
an actual modification is needed.

\section{Creating DataFrames}

\subsection*{\texttt{input} block}

For small inline datasets — ideal for examples and tests:

\begin{lstlisting}[caption={input block}]
input x y z
  1   2  10
  3   4  20
  5   6  30
end
\end{lstlisting}

The \lstinline{input} block creates a DataFrame named \lstinline{df} by
default. The number of columns is defined by the header; each data row
must have the same number of values.

\begin{nota}
  \lstinline{input} accepts only numeric values. Use \lstinline{.} to
  represent missing values. For data with text columns, use
  \lstinline{load}.
\end{nota}

\subsection*{\texttt{load} — external files}

\begin{lstlisting}[caption={Loading files}]
load "data.csv" as sales
load "panel.dta" as panel        // Stata format
load "results.parquet" as df
load "sheet.xlsx" as df          // first sheet
load "sheet.xlsx" sheet="Data" as df
\end{lstlisting}

Supported formats: CSV, Parquet, Stata (\texttt{.dta}), Excel
(\texttt{.xlsx}), SQLite.

\section{Saving DataFrames}

\begin{lstlisting}
save df "output.csv"
save df "output.parquet"
save df "output.dta"
save df "output.xlsx"
\end{lstlisting}

\section{Inspecting DataFrames}

\subsection*{\texttt{display}}

Displays the first rows of the DataFrame:

\begin{lstlisting}
display df
display df, 20    // display up to 20 rows
\end{lstlisting}

\subsection*{\texttt{describe}}

Displays the structure: number of observations, variables, and types:

\begin{lstlisting}
describe df
\end{lstlisting}

Typical output:

\begin{lstlisting}[style=inline]
DataFrame: 1250 obs, 8 vars
  year    Float  (0 missing)
  state   Float  (0 missing)
  income  Float  (12 missing)
  gini    Float  (3 missing)
\end{lstlisting}

\subsection*{\texttt{count} and \texttt{nrow}}

Returns the number of rows:

\begin{lstlisting}
let n  = count(df)               // total observations
let n2 = nrow(df)                // alias for count(df)
let n3 = count(df, income > 1000) // conditional count
\end{lstlisting}

\section{Special variables}

Inside row-by-row commands (\lstinline{generate}, \lstinline{mutate},
\lstinline{filter}), two implicit variables are available:

\begin{description}
  \item[\lstinline{_n}] Current row number (starts at 1).
  \item[\lstinline{_N}] Total rows in the DataFrame.
\end{description}

\begin{lstlisting}
generate df id     = _n       // 1, 2, 3, ..., N
generate df weight = 1.0/_N   // uniform weight
\end{lstlisting}

\section{Time series operators}

Inside \lstinline{generate} expressions (imperative form), special
time-series operators are available:

\begin{center}
\begin{tabular}{cll}
\toprule
\textbf{Operator} & \textbf{Meaning} & \textbf{Example} \\
\midrule
\texttt{L.x}   & Lag (previous value)   & \texttt{L.gdp}  → gdp$_{t-1}$ \\
\texttt{F.x}   & Lead (next value)      & \texttt{F.rate} → rate$_{t+1}$ \\
\texttt{D.x}   & First difference       & \texttt{D.gdp}  → gdp$_t$ - gdp$_{t-1}$ \\
\texttt{L2.x}  & Lag of order 2         & \texttt{L2.gdp} → gdp$_{t-2}$ \\
\bottomrule
\end{tabular}
\end{center}

\begin{lstlisting}[caption={Time series operators}]
generate df dgdp    = D.gdp     // GDP growth
generate df gdp_lag = L.gdp     // lagged GDP
\end{lstlisting}
